\section{Notation}
\label{sec:notation}

This section fixes the symbols used across the specification. It defines only
shared notation and reusable primitive relations; concrete record layouts,
script execution, and signature-message construction are defined in their
respective sections.

\subsection{Types and values}

\begin{longtable}{@{}L{0.24\linewidth}L{0.64\linewidth}@{}}
\toprule
Notation & Meaning \\
\midrule
\(\mathbf{Bool}\) & Truth values \(\{\bot,\top\}\). \\
\(\mathbf{Byte}\) & The integers \(\{0,\ldots,255\}\). \\
\(\mathbf{Bytes}\) & Finite byte strings. \\
\(\mathbf{Bytes}[n]\) & Byte strings of exactly \(n\) bytes. \\
\(\mathbf{N}\), \(\mathbf{Z}\) & Nonnegative integers and integers. \\
\(x:T\) & Value \(x\) has type \(T\). \\
\(f:S \to T\) & Total function from \(S\) to \(T\). \\
\(f(x)=\bot\) & \(f\) is undefined, fails, or rejects input \(x\). \\
\bottomrule
\end{longtable}

\subsection{Sequences and records}

\begin{longtable}{@{}L{0.24\linewidth}L{0.64\linewidth}@{}}
\toprule
Notation & Meaning \\
\midrule
\([x_0,\ldots,x_{n-1}]\) & Ordered sequence with zero-based indexing. \\
\(x[i]\) & Element at index \(i\). \\
\(|x|\) & Length of sequence or byte string \(x\). \\
\(x || y\) & Concatenation of sequences or byte strings. \\
\bottomrule
\end{longtable}

Hex strings are written in display order. Sections that describe serialized
fields state whether display order differs from wire order.

Record tables define field order, type, and meaning. Unless explicitly marked
as logical state, a record table gives the order used by \(\mathrm{E}_{T}\).

\subsection{Encoding and byte order}

For a type or record \(T\), \(\mathrm{E}_{T}(x)\) is the canonical byte
encoding of \(x\) as \(T\). The partial inverse
\(\mathrm{E}_{T}^{-1}(b)\) parses bytes \(b\) as \(T\) and is undefined for
non-canonical or invalid encodings. When parsing succeeds,
\(\mathrm{E}_{T}^{-1}(\mathrm{E}_{T}(x))=x\).

\(\mathrm{le}(b)\) is the unsigned little-endian integer represented by byte
string \(b\). Fixed-width integer encodings and CompactSize are defined in
Section~\ref{sec:serialization}.

\subsection{Hashes}

Hash values are protocol bytes. Human display strings often reverse the bytes
of 256-bit hashes.

\[
\begin{aligned}
  \mathrm{H}(b) &= \operatorname{SHA256}(\operatorname{SHA256}(b)),\\
  \operatorname{HASH160}(b) &= \operatorname{RIPEMD160}(\operatorname{SHA256}(b)),\\
  \operatorname{TagHash}_{tag}(b)
    &= \operatorname{SHA256}(\operatorname{SHA256}(tag) ||
       \operatorname{SHA256}(tag) || b).
\end{aligned}
\]

\(\mathrm{H}\) is \code{hash256} in Bitcoin source material. Taproot and
BIP340 Schnorr signatures use \(\operatorname{TagHash}\)
\cite{nakamoto-bitcoin,bip-0340,bip-0341}.

For a nonempty list of 32-byte hashes \(L\), \(\mathrm{MR}(L)\) is the Merkle
root obtained by repeatedly replacing adjacent pairs \((a,b)\) with
\(\mathrm{H}(a||b)\). If a level has odd length, the final hash is duplicated
for that level. \(\mathrm{MR}([])=\bot\).

\subsection{Identifiers}

\[
\begin{aligned}
  txid(tx) &= \mathrm{H}(\mathrm{E}_{0}(tx)),\\
  wtxid(tx) &=
    \begin{cases}
      \mathrm{H}(\mathrm{E}_{w}(tx)), & \text{witness serialization is used},\\
      txid(tx), & \text{otherwise},
    \end{cases}\\
  \mathrm{bh}(h) &= \mathrm{H}(\mathrm{E}_{\mathrm{BlockHeader}}(h)).
\end{aligned}
\]

\subsection{Chain context}

For a chain \(C\), \(C[k]\) is the block at height \(k\), with \(C[0]\) as
genesis. \(\operatorname{height}(B)\) is the height of block \(B\), and
\(\operatorname{tip}(C)=C[|C|-1]\). \(\operatorname{MTP}(C,k)\) is the median
timestamp of \(C[k]\) and its ten nearest ancestors, or all available ancestors
when \(k<10\).

\subsection{Cryptographic relations}

All public-key signature relations use the secp256k1 group.
\(\operatorname{ECDSAVerify}(Q,m,\sigma)\) is true iff ECDSA signature
\(\sigma\) verifies digest \(m\) under public key \(Q\), subject to active script
encoding rules. \(\operatorname{SchnorrVerify}(Q,m,\sigma)\) is the BIP340
verification relation \cite{bip-0066,bip-0340}.
